% !TeX root = ../教材题库_全部条目_转写工作稿.tex
% Source: XJ-B1-C06-TH-P217-Y1023
% 第6章 / 6.2 抽样 / 教材 p.217 / PDF 第224页
% OCR未校对：公式与文字均可能有误；未完成数学TeX转写；禁止直接发给学生
% 原题及相邻原解截图（核对用）：../题目原图/XJ-B1-C06-TH-P217-Y1023-01.webp
% 本文件只在转写完成后作为题干使用；OCR 条目尚不是可靠题干。
\begin{bankitem}{p.217 思考（含紧邻说明）}{第6章 / 6.2 抽样 / 教材 p.217 / PDF 第224页}{XJ-B1-C06-TH-P217-Y1023}{OCR未校对：公式与文字均可能有误；未完成数学TeX转写；禁止直接发给学生}
\begin{ocrdraft}
[教材p.217 / PDF第224页]
\par
？思考
\par
简单随机抽样是一种最基本的抽样方法，是其他抽样方法的基础．上
\par
述两种常用简单随机抽样方法各有什么特点？
\par
抽签法的突出特点是简单、直观，在总体个数不大时，使总体处于“搅拌均匀”
\par
的状态容易实现，这时，每个个体有均等的机会被抽中，从而能够保证样本的代表
\par
性．但当总体中个体很多时，对个体编号的工作量很大，并且难以搅拌均匀，可能
\par
导致抽样的不公平，采用随机数的方法可以克服这一问题．
\end{ocrdraft}
\end{bankitem}
